\section{DD1：线性正则性下的补偿关系}\label{sec:dd1}

本节在线性残差正则性下给出两个方向的分段补偿关系：由正则性确定允许的次单调性缺陷，或由缺陷反求所需正则性。结合近端步的局部存在性，距离收缩与比例步长界共同给出有限长度与点收敛。

沿用引言中 $H,F,S,d,P_S,r_F$ 的记号。
固定 $\lambda,\ell,R>0$ 及开集 $U\subseteq H$，令
\[
 U_R=\{x\in U:d(x)\le R\},\qquad L=\ell/\lambda,\qquad
 \mathcal J(x)\subseteq J_\lambda(x)=(I+\lambda F)^{-1}(x).
\]
$\mathcal J$ 表示迭代允许选取的近端输出，不要求单值。设有限函数 $\Phi:[0,R]\to\R$ 满足：对每个 $x\in U_R$ 及 $y\in\mathcal J(x)$，
\begin{equation}\label{dd1:assumptions}
 d(y)\le\ell r_F(y),\qquad
 \exists p\in P_S(x):\quad\ip{x-y}{y-p}\ge\Phi(d(x)).
\end{equation}
这里使用完整集合 $F(y)$ 的真残差；最近点 $p$ 可以随 $(x,y)$ 改变。若误差界只在指定输出邻域或残差范围成立，$\mathcal J$ 的全部输出均须满足相应限制。其非空性另行假设。

\begin{lemma}[一步几何估计]\label{dd1:geometry}
对满足 \eqref{dd1:assumptions} 的近端对，记 $t=d(x)$、$a=d(y)$、$s=\norm{x-y}$、$h=\norm{y-p}$。则
\begin{equation}\label{eq:dd1-energy}
 a\le Ls,\qquad a\le h,\qquad
 a^2+s^2\le h^2+s^2\le t^2-2\Phi(t),
\end{equation}
且 $t^2-4\Phi(t)\ge0$，并有
\begin{equation}\label{eq:dd1-step}
 \max\{s,h\}\le B_\Phi(t),\qquad
 B_\Phi(t)=\frac{t+\sqrt{t^2-4\Phi(t)}}2.
\end{equation}
当 $t=0$ 时，$-s^2\ge\Phi(0)$；特别地，$\Phi(0)\ge0$ 可保证 $y=x$。
\end{lemma}

\begin{proof}
由近端包含关系，$r_F(y)\le s/\lambda$，故 $a\le Ls$。又因 $p\in S$，有 $a\le h$，而
\[
 t^2=s^2+h^2+2\ip{x-y}{y-p}\ge s^2+h^2+2\Phi(t).
\]
分别将内积写为 $\ip{x-y}{x-p}-s^2$ 和 $\ip{x-p}{y-p}-h^2$，可得
$s^2-ts+\Phi(t)\le0$ 与 $h^2-th+\Phi(t)\le0$，其较大根即为 $B_\Phi(t)$。若 $t=0$，则 $p=x$，直接得到最后一项结论。
\end{proof}

以上估计保留 $\Phi$ 的符号。为求统一收缩因子，记
\begin{equation}\label{dd1:negative-part}
 D(t)=[-\Phi(t)]_+,\qquad
 \sigma_R=\sup_{0<t\le R}\frac{D(t)}{t^2}.
\end{equation}

\begin{theorem}[DD11：由正则性确定允许缺陷]\label{thm:dd1}
设 \eqref{dd1:assumptions} 成立且 $\sigma_R<\infty$。对 $\sigma\ge0$，令
\begin{align}
 b_\sigma&=\frac{1+\sqrt{1+4\sigma}}2,
 &e_L(\sigma)&=L\sqrt{\frac{1+2\sigma}{1+L^2}},\label{dd1:constants}\\
 m_L(\sigma)&=\min\{b_\sigma,Lb_\sigma,e_L(\sigma)\},
 &\kappa_L(\sigma)&=\min\{Lb_\sigma,e_L(\sigma)\}.\label{dd1:factors}
\end{align}
则每个正距离输入 $x\in U_R$ 及 $y\in\mathcal J(x)$ 均满足
\begin{equation}\label{dd1:forward-estimates}
 d(y)\le m_L(\sigma_R)d(x),\qquad
 \norm{x-y}\le b_{\sigma_R}d(x).
\end{equation}
其中 $m_L(\sigma)<1$ 等价于 $\kappa_L(\sigma)<1$，亦等价于
\begin{equation}\label{dd1:threshold}
 \sigma<T(L),\qquad
 T(L)=\begin{cases}
 (1-L)/L^2,&0<L\le1/2,\\[2pt]
 1/(2L^2),&L\ge1/2.
 \end{cases}
\end{equation}
更具体地，给定 $r\in(0,1)$，令 $r_0(L)=L/\sqrt{1+L^2}$，并定义
\begin{equation}\label{dd1:target-bound}
 \Sigma_L(r)=\begin{cases}
 \dfrac{r(r-L)}{L^2},&0<L<\dfrac12,\quad\dfrac{L}{1-L}<r<1,\\[6pt]
 \dfrac{(1+L^2)r^2-L^2}{2L^2},&r_0(L)\le r<1\text{ 的其余情形}.
 \end{cases}
\end{equation}
若 $r\ge r_0(L)$ 且
\begin{equation}\label{dd1:target-condition}
 \frac{D(t)}{t^2}\le\Sigma_L(r)\qquad(0<t\le R),
\end{equation}
则 $d(y)\le r d(x)$。两支在 $r=L/(1-L)$ 处取值相同。
若另有零输入驻定条件
\begin{equation}\label{dd1:stationarity}
 x\in U_R\cap S,\quad y\in\mathcal J(x)\quad\Longrightarrow\quad y=x,
\end{equation}
则上述距离与步长估计也适用于 $t=0$；$\Phi(0)\ge0$ 是该条件的一个充分条件。
\end{theorem}

\begin{proof}
由 $\Phi(t)\ge-\sigma_Rt^2$ 及引理~\ref{dd1:geometry}，有
$s,h\le b_{\sigma_R}t$、$a\le Ls$，以及
\[
 a^2(1+L^{-2})\le a^2+s^2\le(1+2\sigma_R)t^2.
\]
这给出 $m_L$ 的三项上界。因 $b_\sigma\ge1$，严格收缩当且仅当
$Lb_\sigma<1$ 或 $e_L(\sigma)<1$。后者等价于 $\sigma<1/(2L^2)$；前者要求 $L<1$，并等价于 $\sigma<(1-L)/L^2$。比较两个上界，分界点为 $L=1/2$，故得 \eqref{dd1:threshold}。

对于给定 $r<1$，$e_L(\sigma)\le r$ 等价于
$\sigma\le\Sigma_E(r)=((1+L^2)r^2-L^2)/(2L^2)$；而在 $r\ge L$ 时，$Lb_\sigma\le r$ 等价于 $\sigma\le\Sigma_B(r)=r(r-L)/L^2$。两者之差为
\begin{equation}\label{dd1:branch-difference}
 \Sigma_B(r)-\Sigma_E(r)
 =\frac{((1+L)r-L)((1-L)r-L)}{2L^2}.
\end{equation}
在 $r\ge r_0(L)$ 时，第一因子为正；第二因子的零点落在 $(0,1)$ 内当且仅当 $L<1/2$。取两项适用上界的较大者，便得 \eqref{dd1:target-bound}。逐个 $t$ 代入 $\sigma=D(t)/t^2$ 即得结论。零距离情形由 \eqref{dd1:stationarity} 保证。
\end{proof}

两支对应不同的最近点构形：当 $s\le t+a$ 时，可以由平方距离关系决定界；当 $s>t+a$ 时，$h\ge s-t>a$，输入最近点不再是输出最近点，此时步长根式给出较强估计。

\begin{theorem}[DD12：由允许缺陷确定正则性]\label{thm:dd12}
设 \eqref{dd1:assumptions} 成立，且 $D(t)\le\sigma t^2$，$0<t\le R$，其中 $\sigma\ge0$。给定 $0<r<1$，若
\begin{equation}\label{dd1:inverse-target}
 \frac\ell\lambda\le\Lambda_\sigma(r),\qquad
 \Lambda_\sigma(r)=\begin{cases}
 \dfrac r{\sqrt{1+2\sigma-r^2}},&0\le\sigma\le r(1+r),\\[6pt]
 \dfrac{2r}{1+\sqrt{1+4\sigma}},&\sigma\ge r(1+r),
 \end{cases}
\end{equation}
则正距离输入均满足 $d(y)\le r d(x)$。两支在 $\sigma=r(1+r)$ 处相同。存在某个统一收缩因子的充分条件为
\begin{equation}\label{dd1:inverse-threshold}
 \frac\ell\lambda<\Lambda(\sigma),\qquad
 \Lambda(\sigma)=\begin{cases}
 +\infty,&\sigma=0,\\[2pt]
 1/\sqrt{2\sigma},&0<\sigma\le2,\\[2pt]
 2/(1+\sqrt{1+4\sigma}),&\sigma\ge2.
 \end{cases}
\end{equation}
该条件等价于 $\sigma<T(\ell/\lambda)$。若 \eqref{dd1:stationarity} 成立，距离估计同时适用于零输入。
\end{theorem}

\begin{proof}
因 $r<1\le b_\sigma$，$m_L(\sigma)\le r$ 当且仅当
\[
 L\le\frac r{b_\sigma}\quad\text{或}\quad
 L\le\frac r{\sqrt{1+2\sigma-r^2}}.
\]
由 $\sigma=b_\sigma(b_\sigma-1)$，两分母平方之差为
$1+2\sigma-r^2-b_\sigma^2=(b_\sigma-1)^2-r^2$，故较大上界在 $\sigma=r(1+r)$ 处切换，得到 \eqref{dd1:inverse-target}。将严格收缩的两个条件写为
$L<1/b_\sigma$ 或 $L<1/\sqrt{2\sigma}$，分界为 $\sigma=2$，即得 \eqref{dd1:inverse-threshold}；$\sigma=0$ 时第二条件对所有有限 $L$ 成立。
\end{proof}

\begin{theorem}[有限长度与局部收敛]\label{thm:dd1-trajectory}
设 \eqref{dd1:assumptions}、\eqref{dd1:stationarity} 成立，$\sigma_R<T(L)$，并令
$b=b_{\sigma_R}$、$\kappa=\kappa_L(\sigma_R)$、$C=b/(1-\kappa)$。
每条始终满足 $x^{k+1}\in\mathcal J(x^k)$、$x^k\in U_R$ 的无限序列均有
\begin{align}
 d(x^k)&\le\kappa^k d(x^0),&
 \sum_{j=k}^\infty\norm{x^{j+1}-x^j}&\le C d(x^k),\label{dd1:trajectory-estimates}\\
 x^k&\to x^\infty\in S,&
 \norm{x^k-x^\infty}&\le C\kappa^k d(x^0).\label{dd1:point-rate}
\end{align}
若每个 $x\in U_R$ 都满足 $\mathcal J(x)\ne\varnothing$，且初值 $x^0\in U_R$ 满足
\begin{equation}\label{dd1:localization}
 C d(x^0)<\dist(x^0,U^c),
\end{equation}
则任意逐步选取的迭代都可无限继续，始终位于 $U_R$，且极限属于 $S\cap U$。约定 $\dist(x,\varnothing)=+\infty$。
\end{theorem}

\begin{proof}
记 $t_k=d(x^k)$、$s_k=\norm{x^{k+1}-x^k}$。DD11 给出 $t_{k+1}\le\kappa t_k$ 和 $s_k\le bt_k$。对已经存在的无限序列求和，便有
$\sum_{j=k}^\infty s_j\le bt_k/(1-\kappa)=Ct_k$。序列因此为柯西列，结合 $t_k\to0$ 及 $S$ 闭，得到极限与误差估计。

为证明局部续接，先考虑任意有限迭代。每步非空性提供下一输出，一步估计给出
\[
 d(x^{k+1})\le\kappa^{k+1}t_0\le R,\qquad
 \norm{x^{k+1}-x^0}\le bt_0\sum_{j=0}^k\kappa^j\le Ct_0<\dist(x^0,U^c).
\]
故该输出仍在 $U_R$，归纳即可无限继续。所有迭代点及极限均位于包含在 $U$ 内的闭球 $\overline B(x^0,Ct_0)$ 中，结论随之成立。
\end{proof}

\begin{example}[超出已有兼容条件的完整算子]\label{dd1:example}
取 $H=\R^2$、$\lambda=1$、$\ell=1/4$，令
\[
 b=\frac{1+\sqrt{41}}2,\quad h=b-1,\quad a=b/4,\quad
 c=\sqrt{h^2-a^2},\quad\bar y=(c,a),\quad\bar x=-\bar y/h.
\]
此时 $h>a>0$。设 $S=\R_+e_1$，定义
\begin{equation}\label{dd1:example-operator}
 F(z)=\begin{cases}
 \{(v_1,v_2):v_1\le0\},&z=0,\\
 \{(0,v_2):v_2\in\R\},&z=\xi e_1,\quad\xi>0,\\
 \{\tau(\bar x-\bar y)\},&z=\tau\bar y,\quad\tau>0,\\
 \varnothing,&\text{其余}.
 \end{cases}
\end{equation}
$F$ 有闭图、$\zer F=S$，且直接求逆得到
\begin{equation}\label{dd1:example-resolvent}
 \gph J_1=\gph P_S\cup\{(\tau\bar x,\tau\bar y):\tau\ge0\},
 \qquad P_S(w)=\max\{w_1,0\}e_1.
\end{equation}
因此每个输入的完整近端映射均非空。对额外近端对 $x=\tau\bar x$、$y=\tau\bar y$，最近输入点为 $p=0$，并有
\[
 d(x)=\tau,\quad d(y)=a\tau,\quad s=b\tau,
 \quad r_F(y)=b\tau,\quad\ip{x-y}{y-p}=-10\tau^2.
\]
这里 $\norm{\bar y}=h$、$\norm{\bar x}=1$、$b(b-1)=10$，且 $F(y)$ 仅含定义中的单个向量。其余近端对均为投影对，输出距离、真残差与相应内积都为零。因此全局满足 \eqref{dd1:assumptions}，可取 $\Phi(t)=-10t^2$；系数 $10$ 和 $1/4$ 在任意小邻域内均不能减小。又有
\[
 T(1/4)=12>10,\qquad
 \kappa_{1/4}(10)=\frac{1+\sqrt{41}}8<1,
 \qquad e_{1/4}(10)=\sqrt{\frac{21}{17}}>1.
\]
DD1 因而给出距离收缩。事实上，额外输出的第一坐标为正，不再是额外输入，故所有轨道至多经两次非驻定迭代便进入 $S$。

将文献中的条件换算为相同 $\lambda$ 与输出误差常数，Luke--Tam 的显式条件为 $\sigma(1+L)^2<1/2$，见 \cite[假设~2(d)(ii)、定理~2]{LukeTam2025}；Alcantara 的能量条件为 $2\sigma L^2<1$，见 \cite[注~3.5]{Alcantara2026}。在 $L=1/4$ 时，三者分别要求
\[
 \begin{array}{c|c}
 \text{充分条件}&\text{允许的系数}\\\hline
 \text{Luke--Tam}&\sigma<8/25\\
 \text{Alcantara 的能量条件}&\sigma<8\\
 \text{DD1}&\sigma<12.
 \end{array}
\]
本例满足 DD1 而不满足前两项。DD1 中的能量因子已见于 Alcantara 的 注~3.5，其 命题~3.4 的证明也给出了步长根式估计。将后者与残差界相乘，再与能量因子取较小值，得到 $0<L<1/2$ 时新增的范围
$1/(2L^2)\le\sigma<(1-L)/L^2$。上述比较针对所列兼容条件，完整局部定理的其他假设仍各自保留。
\end{example}
